\chapter{Dettagli Implementativi}
\label{chap:implementazione}

Questo capitolo analizza in dettaglio le soluzioni tecniche e gli algoritmi adottati per implementare le funzionalità chiave del sistema. Verranno presentati e commentati degli estratti di codice significativi, focalizzandosi sulle parti più rilevanti del backend Node.js e della logica custom del client.

\section{Implementazione del Backend Node.js}

\subsection{Gestione delle Richieste e Routing}
Il server è stato implementato utilizzando Express.js, che facilita la definizione delle rotte e la gestione delle richieste HTTP. Di seguito la configurazione base del server:

\begin{verbatim}
const express = require('express');
const session = require('express-session');
const app = express();
const port = 3000;

// Middleware per il parsing del corpo delle richieste
app.use(express.json({ limit: '10mb' }));
app.use(express.urlencoded({ extended: true }));

// Configurazione delle sessioni
app.use(session({
  secret: 'my-secret-key',
  resave: false,
  saveUninitialized: false,
  cookie: { secure: false, maxAge: 24 * 60 * 60 * 1000 } // 24 ore
}));

// Servizio dei file statici
app.use(express.static('.'));

// Definizione delle rotte API
app.put('/data/data.json', isAuthenticated, handleDataSave);
app.get('/api/search-publications', handlePublicationSearch);
app.post('/login', handleLogin);
app.post('/logout', handleLogout);
app.get('/check-auth', handleCheckAuth);
app.post('/create-blank-page', isAuthenticated, handleCreateBlankPage);
app.post('/create-page-from-template', isAuthenticated, handleCreateFromTemplate);

// Avvio del server
app.listen(port, () => {
  console.log(`Server in esecuzione su http://localhost:${port}`);
});
\end{verbatim}

\subsection{Sistema di Autenticazione con Express-session}
La sicurezza delle operazioni di scrittura è garantita da un sistema di sessioni basato sul pacchetto \texttt{express-session}. Un middleware, denominato \texttt{isAuthenticated}, viene anteposto a tutte le rotte che modificano il file system.

\begin{verbatim}
// Middleware di protezione
const isAuthenticated = (req, res, next) => {
  if (req.session.isAuthenticated) {
    return next(); // Utente autenticato, prosegui
  }
  res.status(401).json({ error: 'Accesso negato. Effettuare il login.' });
};

// Gestione del login
app.post('/login', (req, res) => {
  const { username, password } = req.body;
  
  // Verifica semplificata delle credenziali
  if (username === 'admin' && password === 'password') {
    req.session.isAuthenticated = true;
    req.session.username = username;
    res.json({ success: true, message: 'Login effettuato con successo' });
  } else {
    res.status(401).json({ success: false, error: 'Credenziali non valide' });
  }
});

// Gestione del logout
app.post('/logout', isAuthenticated, (req, res) => {
  req.session.destroy((err) => {
    if (err) {
      console.error('Errore durante il logout:', err);
      return res.status(500).json({ error: 'Errore durante il logout' });
    }
    res.clearCookie('connect.sid');
    res.json({ success: true, message: 'Logout effettuato con successo' });
  });
});

// Verifica stato autenticazione
app.get('/check-auth', (req, res) => {
  res.json({ isAuthenticated: !!req.session.isAuthenticated });
});
\end{verbatim}

Questo semplice ma efficace meccanismo assicura che solo gli utenti che hanno completato con successo il login possano accedere alle funzionalità di modifica.

\section{Automatizzazione del Versioning su Git}
Il requisito fondamentale di integrare un sistema di version control è stato risolto automatizzando i comandi Git direttamente dal server Node.js ad ogni operazione di salvataggio.

\begin{verbatim}
const { exec } = require('child_process');
const fs = require('fs').promises;

// Gestione salvataggio dati e commit automatico
app.put('/data/data.json', isAuthenticated, async (req, res) => {
  const filePath = './data/data.json';
  const content = JSON.stringify(req.body, null, 2); // Formattazione leggibile
  
  try {
    // Scrittura del file
    await fs.writeFile(filePath, content, 'utf8');
    console.log('File data.json salvato con successo.');

    // Esecuzione comandi Git
    const timestamp = new Date().toISOString();
    const commitMessage = `Aggiornamento contenuti - ${timestamp}`;
    
    exec(`git add "${filePath}"`, (addErr, addStdout, addStderr) => {
      if (addErr) {
        console.error('Errore git add:', addStderr);
        return res.status(500)
          .json({ error: 'File salvato ma errore in git add' });
      }
      
      exec(`git commit -m "${commitMessage}"`, 
        (commitErr, commitStdout, commitStderr) => {
        if (commitErr) {
          console.error('Errore git commit:', commitStderr);
          return res.status(500).json({ 
            error: 'File salvato ma errore in git commit' 
          });
        }
        
        console.log('Commit Git eseguito:', commitStdout);
        res.json({ 
          success: true, 
          message: 'File salvato e commit eseguito con successo' 
        });
      });
    });
  } catch (error) {
    console.error('Errore scrittura file:', error);
    res.status(500).json({ error: 'Errore durante il salvataggio del file' });
  }
});
\end{verbatim}

Come si evince dal codice, dopo aver scritto con successo il file \texttt{data.json}, viene invocata la funzione \texttt{exec} dal modulo \texttt{child\_process} per eseguire i comandi Git. Questo approccio rende il versioning una parte intrinseca del processo di salvataggio.

\section{Scraping di Dati Esterni con Cheerio}
Per implementare la ricerca di pubblicazioni senza avere accesso a un'API ufficiale, si è ricorso a una tecnica di web scraping lato server.

\begin{verbatim}
const cheerio = require('cheerio');

// Ricerca pubblicazioni da art.torvergata.it
app.get('/api/search-publications', async (req, res) => {
  const { query } = req.query;
  
  if (!query || query.trim().length < 3) {
    return res.status(400).json({ 
      error: 'La query di ricerca deve contenere almeno 3 caratteri' 
    });
  }
  
  try {
    const searchUrl = 
      `https://art.torvergata.it/search?q=${encodeURIComponent(query)}`;
    
    const response = await fetch(searchUrl);
    if (!response.ok) {
      throw new Error(`HTTP error! status: ${response.status}`);
    }
    
    const html = await response.text();
    const $ = cheerio.load(html);

    const publications = [];
    
    // Estrazione dati dalla tabella dei risultati
    $('#tableView_body table tbody tr').each((i, row) => {
      const columns = $(row).find('td');
      if (columns.length >= 3) {
        const titleElement = $(columns[1]).find('a');
        const title = titleElement.text().trim();
        const link = `https://art.torvergata.it${titleElement.attr('href')}`;
        const authors = $(columns[2]).text().trim();
        const year = $(columns[3]).text().trim();
        
        if (title) {
          publications.push({ 
            title, 
            link, 
            authors, 
            year,
            source: 'art.torvergata.it'
          });
        }
      }
    });

    res.json({
      success: true,
      count: publications.length,
      results: publications
    });
  } catch (error) {
    console.error('Errore durante lo scraping:', error);
    res.status(500).json({ 
      error: 'Errore durante la ricerca delle pubblicazioni' 
    });
  }
});
\end{verbatim}

La libreria Cheerio si è rivelata fondamentale: essa fornisce un'API simile a jQuery per navigare e manipolare il DOM in ambiente server. Il codice estrae i dati dalla tabella dei risultati e costruisce un array di oggetti JSON puliti.

\section{Personalizzazione del Client SimplyEdit}
Per integrare le funzionalità del backend con l'interfaccia utente, è stato necessario estendere e personalizzare il comportamento di default di SimplyEdit.

\subsection{Gestione dello Stato di Autenticazione}
\begin{verbatim}
// simply-edit-config.js - Gestione autenticazione
document.addEventListener('DOMContentLoaded', function() {
  // Verifica stato autenticazione
  function checkAuth() {
    return fetch('/check-auth')
      .then(response => response.json())
      .then(data => {
        const isAuthenticated = data.isAuthenticated;
        
        // Mostra/nascondi pulsante modifica
        const editButton = document.getElementById('simply-edit-toggle');
        if (editButton) {
          editButton.style.display = isAuthenticated ? 'block' : 'none';
        }
        
        // Redirect se non autenticato ma si tenta di accedere
        if (!isAuthenticated && window.location.hash === '#simply-edit') {
          window.location.href = '/login.html';
          return false;
        }
        
        return isAuthenticated;
      })
      .catch(error => {
        console.error('Errore verifica autenticazione:', error);
        return false;
      });
  }
  
  // Verifica iniziale
  checkAuth();
  
  // Verifica periodica (ogni 5 minuti)
  setInterval(checkAuth, 5 * 60 * 1000);
});
\end{verbatim}

\subsection{Override delle Funzionalità di Storage}
\begin{verbatim}
// simply-edit-config.js - Override salvataggio pagine
document.addEventListener('simply-storage-init', function() {
  if (!window.editor) return;

  // Sovrascrive la funzione di salvataggio pagine
  editor.storage.page.save = function(url, html) {
    return new Promise((resolve, reject) => {
      let newPagePath = new URL(url, window.location.origin).pathname;
      
      // Determina il tipo di creazione (pagina vuota o da template)
      const isBlankPage = !html;
      const endpoint = isBlankPage ? 
        '/create-blank-page' : '/create-page-from-template';
      
      fetch(endpoint, {
        method: 'POST',
        headers: { 'Content-Type': 'application/json' },
        body: JSON.stringify({ 
          path: newPagePath, 
          html: html 
        })
      })
      .then(response => response.json())
      .then(data => {
        if (data.success) {
          // Reindirizza alla nuova pagina in modalità modifica
          window.location.href = `${newPagePath}#simply-edit`;
          resolve();
        } else {
          console.error('Errore creazione pagina:', data.error);
          reject(new Error(data.error || 'Errore sconosciuto'));
        }
      })
      .catch(error => {
        console.error('Errore richiesta creazione pagina:', error);
        reject(error);
      });
    });
  };
});
\end{verbatim}

Questo dimostra come l'architettura flessibile di SimplyEdit abbia permesso di integrare le sue funzionalità di base con la logica custom del nostro backend.